% Part: turing-machines
% Chapter: machines-computations
% Section: turing-machines

\documentclass[../../../include/open-logic-section]{subfiles}

\begin{document}

\olfileid{tur}{mac}{tur}
\olsection{ટ્યુરિંગ મશીનો}

\begin{explain}
ટ્યુરિંગ મશીનની ઔપચારિક વ્યાખ્યા અમૂર્ત લાગે છે, પણ ખરેખર
સરળ છે: ટ્યુરિંગ મશીન કેવી રીતે કામ કરે છે તે નિર્દિષ્ટ કરવા
જરૂરી બધી માહિતી તે એક જ ગણિતીય રચનામાં સમાવે છે. તેમાં
(1) મશીન કઈ અવસ્થાઓમાં હોઈ શકે, (2) ટેપ પર કયાં પ્રતીકોને
મંજૂરી છે, (3) મશીન કઈ અવસ્થામાં શરૂ થવું જોઈએ અને
(4) મશીનનો સૂચનાઓનો ગણ શું છે, એ બધું આવે છે.
\end{explain}

\begin{defn}[ટ્યુરિંગ મશીન]
\emph{ટ્યુરિંગ મશીન} $M$ એ નીચેની વસ્તુઓનું બનેલું ચતુષ્ટય
$\langle Q, \Sigma, q_0, \delta\rangle$ છે:
\begin{enumerate}
\item \emph{અવસ્થાઓ}નો સાન્ત ગણ~$Q$,
\item $\TMendtape$ અને $\TMblank$ ધરાવતી સાન્ત \emph{વર્ણમાળા}
  $\Sigma$,
\item \emph{પ્રારંભિક અવસ્થા}~$q_0 \in Q$,
\item સાન્ત \emph{સૂચનાઓનો ગણ}~$\delta\colon Q \times \Sigma
  \pto Q \times \Sigma \times \{\TMleft, \TMright, \TMstay\}$.
\end{enumerate}
આંશિક વિધેય~$\delta$ને $M$નું \emph{સંક્રમણ વિધેય}
પણ કહે છે.
\end{defn}

\begin{explain}
આપણે ધારીએ છીએ કે ટેપ માત્ર એક દિશામાં અનંત છે.
તેથી ટેપના ડાબા છેડાને ચિહ્નિત કરવા માટે વિશિષ્ટ
પ્રતીક~$\TMendtape$ રાખવું ઉપયોગી છે. તેનાથી ટ્યુરિંગ
મશીનના પ્રોગ્રામને ટેપની બહાર નીકળી જવાનો ``ભય'' ક્યારે
છે તે ઓળખવાનું સરળ બને છે. આપણે એમ ધારી શકીએ કે આ
પ્રતીક ક્યારેય બદલવામાં આવતું નથી, એટલે
$\delta(q,\TMendtape)$ વ્યાખ્યાયિત હોય તો
$\delta(q,\TMendtape) = \tuple{q', \TMendtape, x}$.
કેટલાંક પાઠ્યપુસ્તકો આવી ધારણા કરે છે; આપણે કરતા નથી.
તમે તમારું ટ્યુરિંગ મશીન ઘડતી વખતે તે ક્યારેય
$\TMendtape$ને બદલે નહીં તેની કાળજી રાખી શકો છો.
ઉપરાંત, કેટલીક પરિસ્થિતિમાં આ પ્રતીક બદલવાની છૂટ
આપવાથી ઉપયોગી લવચીકતા મળે છે.
\end{explain}

\begin{ex}
\emph{સમ મશીન:} ઔપચારિક રીતે સમ મશીન એ ચતુષ્ટય
$\tuple{Q, \Sigma, q_0, \delta}$ છે, જ્યાં
\begin{align*}
Q & = \{ q_0, q_1 \} \\
\Sigma & = \{ \TMendtape, \TMblank, \TMstroke \}, \\
\delta(q_0, \TMstroke) & = \tuple{q_1, \TMstroke, \TMright},\\
\delta(q_1, \TMstroke) & = \tuple{q_0, \TMstroke, \TMright},\\
\delta(q_1, \TMblank)  & = \tuple{q_1, \TMblank, \TMright}.
\end{align*}
\end{ex}

\end{document}
